\section{Introduction}
\label{sec:intro}

Suppose the fully etched strips of a silicon photonic process come out with walls tilted by \qty{1}{\degree} and the
monitor rings are fitted with vertical walls. In our simulations the fit absorbs the tilt as a width change of
\qty{-3.8403}{nm}, and a ring made of a \qty{500}{nm} rib on a \qty{90}{nm} slab, predicted from that width, lands
\qty{710.8}{pm} away from where it actually resonates. For comparison, Zortman \textit{et~al.}\ held microdisk resonances to within
\qty{105}{GHz} inside a die~\cite{zortman2010}, which is \qty{841.5}{pm} at \qty{1550}{nm}. Post-fabrication trimming
sets rings to within \qty{\pm 32}{pm} of target~\cite{jayatilleka2021}. An unmonitored degree of sidewall angle is therefore
an error of the same size as the within-die spread that process monitoring is meant to explain. It arises because a
rib's resonance responds to wall tilt differently from a strip's, so a width that matches the strips is wrong for
the ribs.

The optical test structures that foundries and designers use measure the effective and group indices of a few
waveguides and convert them into a width and a thickness~\cite{xing2018,zortman2010,lu2017}. The sidewall angle
enters that conversion as a fixed input. Xing \textit{et~al.}\ took the \qty{85}{\degree} of the foundry handbook and
noted that the error from the unknown wall shape is very hard to compensate, so that some extracted parameters are
only relative~\cite{xing2018}. In this paper we show that, for oxide-clad strips, the omission is forced by the
structures. To first order the angle cannot be separated from the width for any number of strips, widths,
polarizations or readouts. We then work out which additional structures make the angle identifiable, how much prior
process knowledge they need, and what precision they reach under a measurement model built from published ring
metrology~\cite{janz2023}. We then ask what the practical limit is. The result depends on the assumption, common to all multi-device
extraction, that the structures share one set of process parameters, and most of the paper's practical conclusions concern what
happens when they do not.

This is a computational study of test-structure design and identifiability. Every number comes from full-vector
finite-element mode solutions, Fisher information or nonlinear fits of simulated data, and no device was fabricated.
The strip result is first-order and holds for the symmetric oxide-clad cross-section under the measurement model of
Section~\ref{sec:theory}.

\subsection{Related work}

\textit{Optical extraction of waveguide geometry.} Xing \textit{et~al.}\ fitted the effective and group indices of a
low-order and a high-order MZI, mapped them to linewidth and thickness with a third-order polynomial, and reached
total errors of \qty{0.37}{nm} and \qty{0.26}{nm} on 44 sites of an MPW die; they also showed how the index range
must be bounded for the inversion to pick the right interference order~\cite{xing2018}. Later work extracted the same two dimensions
from a single high-order MZI~\cite{liu2025}, from racetrack resonators feeding a layout-aware variability
model~\cite{lu2017}, from ring biosensors read in TE and TM~\cite{atsumi2012}, and from microdisks read in both
polarizations~\cite{zortman2010}. MZI sets have mapped the effective index and its thermo-optic coefficient across
wafers~\cite{dwivedi2015}. For silicon nitride, unbalanced MZIs determined thickness, width offset and both material
indices across a \qty{300}{mm} wafer~\cite{butt2024}. Rib waveguides have been characterized with two racetracks of
different perimeter, with a cross-section estimate checked against STEM~\cite{zhang2024}. A recent patent
application determines the geometry from many resonances of one ring. In its background, it states that methods
using several devices must assume that the devices share the same width and thickness variations, and that this
condition ``is recognized as often not satisfied''~\cite{nrc2026}. None of these works extracts the sidewall angle.

\textit{Fabrication variability.} Width and thickness vary systematically with layout and position. Hierarchical
models separate layout-dependent from random contributions and tie die-scale linewidth variation to pattern
density~\cite{xing2023}, and the resonance mismatch between identical rings grows linearly with their
separation~\cite{chrostowski2014}. CMOS-line fabrication has reached linewidth control better than \qty{0.6}{nm}
within a chip and \qty{2}{nm} chip-to-chip~\cite{selvaraja2010}. Wafer-scale resonator statistics~\cite{robinson2025,mdpi2025,aim2023}
and variability-aware design flows~\cite{bogaerts2018} show why such numbers matter for yield. Ring basics,
including the relation between resonance, free spectral range and group index, are reviewed
in~\cite{bogaerts2012,chrostowski2015}.

\textit{Sidewall angle and physical metrology.} Angled sidewalls break the vertical symmetry of a waveguide and cause
TM-to-TE hybridization and conversion in tapers~\cite{dai2015,mdpi2021}, the mechanism behind some polarization
rotators~\cite{dai2011}. Sidewall roughness causes backscattering and polarization rotation~\cite{morichetti2010}
and sets the loss floor of high-Q rings~\cite{biberman2012}. Dimensional metrology measures the profile directly.
Cross-sectional TEM is destructive and AFM and SEM are local~\cite{orji2018}. 3D-AFM resolves sidewalls of vertical and
undercut structures with die-level reproducibility better than \qty{2}{\percent}~\cite{yoo2022}. Scatterometry is
fast but model-based, and its inverse problem can have correlated parameters and non-unique
solutions~\cite{raymond2005,orji2018}.

\textit{Experiment design and identifiability.} The Cram\'er--Rao lower bound (CRLB) and Fisher information are the
standard tools for asking how well a parameter can be estimated from a measurement~\cite{kay1993}, and optimal
experiment design chooses measurements to improve them~\cite{pukelsheim2006}. A parameter is structurally
unidentifiable when the model output cannot distinguish it from others for any data~\cite{bellman1970}, and
practically unidentifiable when finite, noisy data leave it unbounded~\cite{raue2009}. In dimensional metrology the
Fisher matrix has been used to choose the wavelength and angles of EUV scatterometry so as to break parameter
correlations~\cite{klein2025}, and simulation studies have mapped the sensitivity limits of optical CD
metrology~\cite{silver2007}. Offset and model errors in inverse scatterometry have been propagated to the fitted
profile~\cite{alassaad2007} and estimated jointly with it~\cite{henn2012}. In optics more broadly, Fisher
information has guided the design of maximally informative illumination~\cite{bouchet2021}.

\subsection{Contribution}

Previous optical extraction leaves the sidewall angle out. Previous identifiability-driven metrology design chooses how
to measure one fixed structure. Here we choose which waveguides to measure, and we show that the choice decides
whether the angle is identifiable at all. Specifically:
\begin{itemize}
\item We show, to first order and confirm with mode solutions, that oxide-clad strips are exactly blind to the
sidewall angle, with the angle/width ratio $R=(y_c-y_\mathrm{pivot})/(h/2)$ set by the centroid of the
wall-perturbation weight.
\item We give a counting argument, confirmed by the Fisher matrix and by nonlinear fits, that one rib from a separate
etch leaves the strip angle unidentifiable, and that a known etch depth does not change this.
\item We find that two widths of the same rib restore identifiability. A ring-only set of five strips and two ribs
bounds the strip angle at \qty{0.393}{\degree} at a \qty{1}{pm} precision, and we test this bound against correlated
noise, calibration offsets, temperature and the readout. Nonlinear fits of data from direct femwell solves reach 1.037
times the bound.
\item We show that fabrication consistency, not readout precision, is the practical limit. Departures of a few
hundredths of a nanometre or degree between structures double the bound. At reported levels of local variation one copy of the
set gives several degrees. We give the replication rule that recovers the bound for random offsets and its dependence on spatial
correlation, show which systematic offsets can be fitted and which make the angle unidentifiable, and give the layouts that keep
process gradients out of the fit and the wavelength reference the set needs.
\end{itemize}

An exhaustive search also finds larger sets with lower Cram\'er--Rao bounds. We report the best set of eight as a bound, beside its
result on femwell data, and do not build the proposal on it.
